\section{A multiscale obstruction for interval compilers}
\label{sec:geometry}

We next examine a restricted converse to the bit-to-word reduction. The
operations below include the endpoints of actual canonical word-transpose
calls. Their internal computations may be adaptive: only their resulting
whole-entry permutations enter the argument. The restriction is on the
surrounding compiler and its public layout, not on how an oracle implements
its endpoint. This complements the local coding and tape-geometry barriers
in \cite{FirstPaper2026,SecondPaper2026}.

\begin{definition}[Permitted stage]
\label{def:geometry-stage}
A bank has addresses $0,\ldots,P-1$. A permitted stage is either a stable
split into even and odd addresses, its inverse, or the following operation
on any public family of disjoint contiguous intervals. Each interval
independently carries one of these operations:
\begin{enumerate}[label=(\roman*),nosep]
\item an arbitrary permutation of $n$ consecutive $k$-bit entries,
      preserving the offset within each entry, where $k,n\geq1$ and
      $kn\leq V(k):=k4^k$;
\item a cyclic rotation by any fixed offset;
\item reversal.
\end{enumerate}
The intervals may have different widths, lengths and starting addresses.
Uncovered addresses are fixed. The parity split requires even $P$.
All choices are independent of the bank contents.
\end{definition}

The formalization constructs the endpoint permutation and its inverse from
these interval and operation data; it does not assume the geometric
conclusion below. A canonical word block of side $r$ and width $b(r)$ is
included because $r\leq2^{b(r)}$, hence $b(r)r^2\leq V(b(r))$.
\lean{PartitionData.toStage} and
\lean{canonicalWordBlock_admissible} establish these two facts.

Fix a positive integer $B$ and partition the bank into aligned bins of
$B$ addresses. An entry-permutation interval of width $k$ is \emph{bad at
scale $B$} when
\[
                     k<B<V(k).
\]
Its whole interval is then marked bad. Rotations, reversals and parity
operations have no bad addresses.

\begin{lemma}[One-stage geometry and exceptional volume]
\label{lem:stage-geometry}
For any permitted stage, the good addresses in one $B$-bin have images in
at most $11$ $B$-bins. The same assertion holds for the inverse stage.
If positive scales $B_0,\ldots,B_{m-1}$ satisfy
$V(B_i)\leq B_j$ whenever $i<j$, the sum of the stage's bad volumes over
all these scales is at most $P$.
\end{lemma}

\begin{proof}[Proof outline]
A good entry block either has total length at most $B$, or has entry width
at least $B$. The small blocks meeting one source bin keep their outputs
within three neighboring bins. At most two disjoint blocks of length
greater than $B$ meet that source bin. In a large entry block, the source
bin meets at most two entries, each translated to at most two output
bins. A rotation splits into at most two translated pieces; a reversal
has at most two output bins. Including the uncovered addresses, these estimates give
$3+2\cdot4=11$. Stable split and interleave meet at most four and six
bins, respectively, and therefore satisfy the same bound. All estimates apply in both directions.

A fixed width $k$ can be bad at at most one separated scale. Summing the
volumes of the disjoint intervals therefore charges each address at most
once. The Lean results are
\lean{Stage.good_image_bins_le}, \lean{Stage.predecessors_card_le}, and
\lean{Stage.total_bad_volume_le}.
\end{proof}

A \emph{path} is a chronological composition of permitted stages. It may
depend on the matrix dimension. Let $\gamma_1,\ldots,\gamma_K$ be paths
of length at most $L$ on $P=d^2$ bits, and let $T$ be the row-major bit
transpose permutation. Define the set of covered source addresses by
\[
 C=\{u:\gamma_j(u)=T(u)\text{ for some }j\}.
\]
Coverage is an endpoint coincidence, not a probability that the compiler
is correct on a random input.

\begin{theorem}[Common-scale path coverage]
\label{thm:path-coverage}
Let $m\geq1$ and let $B_i$ be positive divisors of $d$, with
$B_i\geq B_*>0$ and $V(B_i)\leq B_j$ for $i<j$. Then
\begin{equation}
 \label{eq:path-coverage}
 \frac{|C|}{P}\leq\frac{KL}{m}+\frac{K11^L}{B_*}
 \qquad(d\geq1).
\end{equation}
\end{theorem}

\begin{proof}[Proof outline]
At a fixed scale, mark an original address dirty if its path visits a
bad interval. Since each preceding path prefix is a permutation, pulling
bad sets back preserves their cardinalities. Summing over all scales and
all $K$ paths costs at most $KLP$. Thus one scale has total dirty volume
at most $KLP/m$.

The clean part of one source bin reaches at most $11^L$ bins along each
path. Because $B\mid d$, the transpose sends the $B$ addresses of one
source bin into distinct output bins. A path therefore agrees with the
transpose on at most $11^L$ clean addresses per source bin. There are
$P/B$ source bins. Adding dirty addresses and taking the union over paths
gives \eqref{eq:path-coverage}. The compiled endpoint is
\lean{Compiler.transpose_coverage_bound}; it states the equivalent
unnormalized inequality and also covers the empty bank.
\end{proof}

\begin{corollary}[Pointwise nonlinear readout]
\label{cor:path-readout}
For every fixed $K,L$, there is a positive dyadic integer $d_0$ such
that, for every positive dimension $d$ divisible by $d_0$, no $K$ paths
of length at most $L$, followed by arbitrary public address-dependent
Boolean tables, transpose every bit input. In particular, this excludes
every dyadic $d\geq d_0$.
Specifically, an output of the form
\[
 y_a=f_a\bigl(x_{\gamma_1^{-1}(a)},\ldots,
             x_{\gamma_K^{-1}(a)}\bigr)
\]
cannot equal $x_{T^{-1}(a)}$ for every $x$ and $a$.
\end{corollary}

\begin{proof}[Proof outline]
Choose $m\geq\max\{1,4KL\}$ and a dyadic
$B_*\geq\max\{1,4K11^L\}$, then set $B_{i+1}=V(B_i)$.
Set $d_0=B_{m-1}$. All chosen scales are dyadic, and each divides
$d_0$, hence also every multiple $d$ of $d_0$. The coverage fraction is
at most $1/2$. At an uncovered address,
flipping the required source bit changes the desired output but none of
the table's arguments. This contradicts correctness on the full input
cube. \lean{Compiler.no_bounded_path_transpose_of_dvd} gives this
statement with an explicit $d_0$;
\lean{Compiler.no_bounded_path_transpose} derives the dyadic case.
\end{proof}

The formal circuit model also permits nonlinear gates between stages.
A depth-$L$, fan-in-$q$ node applies a permitted stage to each child's
output and then an arbitrary address-dependent Boolean table. Leaves
apply public unary tables to the original input, allowing constants as
well as copies. The exact recursive unrolling has $q^L$ paths, each of
length $L$. Consequently the same exclusion holds for this circuit
model, at every positive multiple of the corresponding threshold for
$K=q^L$. This follows from \lean{Circuit.unroll} and
\lean{Circuit.no_bounded_circuit_transpose_of_dvd}. The unrolling is
proved from the circuit evaluation rule, rather than supplied as a
dependency assumption.

The theorem does not allow the input to select the interval family or
the paths. Nor does it cover arbitrary cross-address computation,
increasing depth or fan-in, or an unrestricted tape implementation.
It therefore excludes a compiler class, not the unrestricted converse
$\WORD\Rightarrow\BIT$.
